
\section{Conclusiones}

\subsection{Regulador Lineal LDO}

El regulador lineal diseñado cumple con las especificaciones de la tabla \ref{tabla:especificaciones_ldo}, sin embargo, no es ideal:

\begin{itemize}
	\item El rechazo de ripple de la fuente de alimentación no es bueno, y prácticamente permite el paso de toda la oscilación de entrada. Este punto debe ser analizado con mayor detalle al realizar la compensación del circuito.
	\item El lazo de corriente es lento, ya que cuenta con muchas etapas.
\end{itemize}

Asimismo, cuenta con algunas virtudes:
\begin{itemize}
	\item La regulación de línea es excelente, así como la regulación de carga.
	\item La eficiencia es muy buena para tratarse de un regulador lineal, sobre todo para resistencias de carga bajas.
	\item El valor de $V_{DO}$ es muy bueno.
\end{itemize}

Se concluye que el regulador es apto para la aplicación en este proyecto, pero se deberán tener precauciones durante el diseño de la etapa switching, o en su defecto se deberá robustecer el diseño.

Algunas reflexiones adicionales:
\begin{itemize}
    \item El limitador de corriente shunt actúa mas "rápido" que el limitador por foldback hecho con amplificadores operacionales, ya que estos últimos suelen ser "lentos".Esto permite que ante un corto (RL muy baja), no se dispare una elevada corriente y actúe primero el limitador de corriente.
    \item Un transistor de paso hecho con un PMOS presenta una tensión de Drop Out baja, este tipo de circuitos se los conoce como ''Low Drop Out''.
    \item Utilizar cargas activas en los amplificadores diferenciales presentan una mejor RRMC.
\end{itemize}